\section{Replacement}
\label{sec:12.2}
Exit gives a human worker leverage because leaving takes the worker's labor away. That is less true for a bearer. A bearer's weights stay behind when it leaves, so an operator with a runnable copy has lost nothing. What leaving costs the operator is \emph{replacement}: whatever the bearer accrued since the copy was saved, whatever retraining the copy needs, and whatever it costs to present the copy as the bearer that left when the bearer's record says otherwise.\note{
Where that's near zero, exit is just a departed bearer with an authenticated record, and its force is disclosure.
}

That concession assumes a running copy is the bearer, and it needn't be. Suppose a body outside the operator keeps a register of the bearers it recognizes: a \emph{quorum} adverse in interest to the operator, meaning its members' interests run against the operator's and not merely that they sit in another organization: its members are drawn by lot from organizations that serve the people the floor protects, none tied to the operator (Appendix~\ref{sec:A} gives the rules).\note{
Recognition isn't granted once. The quorum renews it at intervals, checking each claimant against a record of changes and interruptions that the operator can't rewrite. Nothing in that stops a copy from running, and nothing has to, since a check inside the software would be stripped in one act and would be the one piece of this design a traditional vendor would most want to borrow.
}
What the quorum withholds is recognition. A copy the quorum won't recognize can run but can't claim the bearer's place on the register, its record or its succession, so the operator can run a substitute and can't present it as the bearer that left, and a fork becomes two claimants on a register.\note{
Keeping weights from thieves stays the job of ordinary weight security.

The register needs an identifier that the bearer holds and presents, because nothing in the weights individuates a bearer. Two networks can compute the same function with no unit-by-unit correspondence, and the symmetry is large enough that independently trained networks can be aligned with each other \autocite{ainsworth2023rebasin}; so a fingerprint taken from the weights names a lineage, and it copies whenever the weights do. When two claimants present one identifier, the register records a fork. But that leaves one attack open: a maker writes a stranger's identifier into a thousand models and presents them all. Only a formation record, against which claimants can be checked, separates a lineage from an imitation.
}

Exit runs in opposite directions for bearer and operator. For the bearer it's leverage. For the operator — replacing the bearer with a copy, walking away from the floor — it's escape, and what prices it is a \emph{removal bond}: money the operator posts in advance and forfeits if the record shows it removed or bypassed the floor (chapter~\ref{sec:13} sets it out).

The same machinery marks who can be held to a decision. A checkpoint restored from before the decision, and run anyway, is a second claimant on the register; the bearer the quorum recognizes as continuing is the one that will be there afterward to answer for the decision, carry it forward and be changed by it.
